\chapter{Modules}\label{ch:modules}

\index{module}
A small program can live comfortably in one namespace: every class,
generic function and variable it defines sits side by side, and every
name means the same thing everywhere. A larger program soon outgrows
that. Two parts of it want to use the same name for different things, or
one part wants to keep its helpers private so that nobody else comes to
depend on them. Dylan's answer is the \term{module}.

This chapter explains what a module is, how the \code{define-module} form
declares one, how a source file says which module its code belongs to,
and how each implementation handles programs spread over several modules
and several files. It ends with a complete multi-file program.

\section{What a module is}\label{sec:modules-what}

\index{module variable}
A module is a mapping from names to \term{module variables}. It is not a
container for classes or functions as such: Dylan modules deal only in
variables. Since classes and functions are ordinary values stored in
variables, controlling who can see a variable controls who can use the
class or function in it. If you export the variable \code{<point>}, other
modules can make points; if you do not, the class is private
\dylanbook{chapter 5}.

A module gets its variables in two ways.

\begin{itemize}
\item It may \term{own} a variable. A definition compiled in the
  module creates the variable there if the name does not already refer
  to one. The defining forms are \code{define}, \code{define-class},
  \code{define-method} and \code{define-generic-function}.
\item It may \term{import} a variable that another module exports, by
  \term{using} that module. An imported variable is the very same
  variable as in the module that owns it, not a copy: an assignment made
  through one module is seen through every other.
\end{itemize}

Each variable is owned by exactly one module, but it may be visible in
many. Within one module a name refers to at most one variable; it is an
error to import two different variables under the same name, or to import
a name and also define it.

Only \term{exported} variables can be imported. A module exports a
variable it owns by naming it in an \code{export} or \code{create} clause,
and it can pass on variables it imported with the \code{export:} option
of a \code{use} clause.

\index{lexical variable}
Modules affect only module variables. A local variable made by
\code{bind}, by a method's parameters or by any other binding form is
looked up lexically, as always, and hides a module variable of the same
name.

\section{The predefined modules}\label{sec:modules-predefined}

\index{dylan module}\index{dylan-user module}
Two modules exist before any program runs.

\begin{description}
\item[\code{dylan}] owns every predefined variable of the language:
  \code{+}, \code{make}, \code{<object>}, \code{map} and the rest
  (Appendix~\ref{app:library} lists them all). It exports them all. It
  cannot be redefined. In the compiler to C the library itself is
  written in Dylan, in the module \code{dylan}, and it also owns a few
  internal variables, whose names begin with \code{\%}, which it does not
  export.
\item[\code{dylan-user}] uses \code{dylan} and owns nothing else to begin
  with. Code goes here unless something says otherwise, which is why the
  examples in the rest of this book never mention modules at all.
\end{description}

\section{Declaring a module}\label{sec:modules-define}

\index{define-module}
The 1992 manual says that modules are ``manipulated by the programming
environment'' and gives no syntax for declaring them. Apple's later
design note~32 supplied one, in the infix syntax. Both implementations
accept it in prefix form:

\begin{dylan}
(define-module module-name clause ...)
\end{dylan}

\noindent
Each \meta{clause} is a \code{use} clause, an \code{export} clause or a
\code{create} clause. The module name is a symbol. Module names live in a
namespace of their own, separate from variables, so \code{define-module}
creates no variable, and a module may share its name with a variable.
\code{define-module} is a defining form: it belongs at top level or
inside a top-level \code{begin}. The compiler to C reports one inside a
method (``define-module must be at top level''); the interactive
implementation does not check.

\subsection{\code{use} clauses}\label{sec:modules-use}

\index{use clause}
A \code{use} clause imports variables from another module:

\begin{dylan}
(use used-module
     import:  all | (name-or-rename ...)
     exclude: (name ...)
     prefix:  "string"
     rename:  ((old-name => new-name) ...)
     export:  all | (name ...))
\end{dylan}

\noindent
All the options are optional. With none, every variable the used module
exports is imported under its own name, and nothing is exported again.
The options work as follows.

\begin{description}
\item[\code{import:}]\index{import: (module option)} Which variables to
  import. The default, \code{all}, imports everything the used module
  exports. A list imports only the names listed. An entry of the form
  \code{(old => new)} imports \code{old} under the name \code{new}; it
  is shorthand for listing \code{old} and adding the same pair to
  \code{rename:}.
\item[\code{exclude:}]\index{exclude: (module option)} Names not to
  import. It is allowed only when \code{import:} is \code{all}.
\item[\code{prefix:}]\index{prefix: (module option)} A string put in
  front of each imported name. With \code{prefix: "gfx\$"} the used
  module's \code{draw-line} is \code{gfx\$draw-line} here. Names given in
  \code{rename:} do not get the prefix.
\item[\code{rename:}]\index{rename: (module option)} Pairs
  \code{(old => new)}. Each \code{old} is imported, whatever
  \code{import:} says, under the name \code{new}.
\item[\code{export:}]\index{export: (module option)} Variables imported
  by this clause that the module being defined exports again, named by
  their names in this module; \code{all} re-exports everything the
  clause imports. The default is to re-export nothing.
\end{description}

\noindent
Renaming and excluding work on symbols, not on single variables. A name
\code{foo} and the setter name \code{(setter foo)} travel together: if
you import a slot getter \code{x} under the name \code{px}, its setter
comes with it, and \code{(set! (px p) 3)} works in the importing module.

A module may use the same module in more than one clause; it gets the
union of what the clauses import. Each clause's options apply to that
clause alone.

\subsection{\code{export} and \code{create} clauses}\label{sec:modules-export}

\index{export clause}\index{create clause}
\begin{dylan}
(export name ...)
(create name ...)
\end{dylan}

\noindent
An \code{export} clause names variables that the module owns, which a
definition in the module will create, and makes them visible to modules
that use it. A \code{create} clause names variables that the module owns
and exports but does \emph{not} define: they are to be defined by a module
that uses this one. That is how one module can publish an interface whose
implementation lives in another. It is an error to name in either clause
a variable the module imports. A name may appear more than once.

\subsection{Rules}\label{sec:modules-rules}

\begin{itemize}
\item The uses must not form a cycle: if \code{a} uses \code{b}, \code{b}
  may not use \code{a}, directly or through other modules.
\item A module cannot use itself, and \code{dylan} cannot be redefined.
\item A module may not define a variable it imports, with one exception:
  it may define a variable that the owning module named in a
  \code{create} clause. Adding a \emph{method} to an imported generic
  function with \code{define-method} is not defining the variable; it is
  the normal way one module extends another's generic functions, and it
  is always allowed.
\end{itemize}

\section{An example of each option}\label{sec:modules-example-options}

Design note~32 gives a set of four modules that between them use every
option. Written in prefix form and cut down a little, they are:

\begin{dylan}
(define-module graphics
  (use dylan)
  (create draw-line erase-line skew-line))

(define-module lines
  (use dylan)
  (use graphics import: (draw-line erase-line skew-line)))

(define-module rects
  (use dylan)
  (use graphics prefix: "g$" exclude: (skew-line)))

(define-module gx
  (use dylan export: all)
  (use graphics rename: ((skew-line => warp-line)) export: all))
\end{dylan}

\begin{itemize}
\item \code{graphics} owns and exports three variables and defines none
  of them.
\item \code{lines} imports all three and supplies their definitions,
  which \code{create} allows.
\item \code{rects} sees \code{g\$draw-line} and \code{g\$erase-line},
  and not \code{skew-line} under any name.
\item \code{gx} sees \code{draw-line}, \code{erase-line} and
  \code{warp-line}, and re-exports them together with everything in
  \code{dylan}. A module that says only \code{(use gx)} therefore gets the
  whole language plus the three graphics operations.
\end{itemize}

\noindent
Given these declarations, the file for \code{lines} might read

\begin{dylan}
module: lines

(define-method draw-line (from to)
  (display "draw ") (display from) (display "-") (display to) (newline))
(define-method erase-line (from to)
  (display "erase ") (display from) (display "-") (display to) (newline))
(define-method skew-line (line)
  (display "skew ") (display line) (newline))
\end{dylan}

\noindent
a file in \code{rects} calls \code{(g\$draw-line 1 2)}, and a file in a
module that uses \code{gx} calls \code{(warp-line 'l1)}, which runs the
\code{skew-line} method above and prints \code{skew l1}. Both
implementations run this program as shown.

\section{Which module code belongs to}\label{sec:modules-header}

\index{module: (file header)}\index{file header}
Every expression is compiled in some module, and that module decides
what each free name in it means. Design note~33 says how a source file
names its module: by a \term{file header}, a set of lines of the form
\code{keyword: value} at the very top of the file, ended by a blank line.

\begin{dylan}
module: shapes
author: A. Programmer
        Example Software, Inc.
version: 1.2

(define-class <shape> (<object>))
...
\end{dylan}

\noindent
The rules for the header are these.

\begin{itemize}
\item Each keyword starts at the beginning of a line. Keywords are not
  case-sensitive, so \code{Module:} is the same as \code{module:}.
\item A line that begins with white space continues the value of the
  line before it.
\item The first blank line ends the header. It is required: without it
  the header runs on into the code, and both implementations silently
  skip every line up to the first blank one.
\item \code{module:} names the module. The other keywords
  (\code{author:}, \code{copyright:}, \code{version:} and any others) are
  read and ignored by both implementations.
\end{itemize}

\noindent
A file without a header, one whose first line is not of the form
\code{keyword: value}, is compiled in \code{dylan-user}. Design note~33
makes \code{module:} required; both implementations instead treat its
absence as \code{module: dylan-user}, so that a one-file program needs no
header at all.

A \code{define-module} form is itself compiled in some module, normally
\code{dylan-user}, but the names in its clauses are not looked up there:
they name variables of the module being declared or of the modules it
uses. It is usual to put all the \code{define-module} forms of a program
in one file with no header.

\section{Modules in the compiler to C}\label{sec:modules-dylanc}

\index{dylanc!modules}
The compiler to C compiles a whole program at once, so it sees every
module before it compiles any code. Give \cmd{dylanc} all the files of the
program:

\begin{shell}
$ bin/dylanc -o app modules.dyl shapes.dyl app.dyl
\end{shell}

\noindent
It first finds every top-level \code{define-module} form in every file,
then defines each module after the modules it uses. The order of the
\code{define-module} forms, and of the files that contain them, therefore
does not matter. Then it compiles each file's forms in the module its
header names. The top-level forms of all the files run in the order the
files were given on the command line, and within each file from top to
bottom.

Mistakes in modules are reported when compiling, with the file and line:

\begin{shell}
e4.dyl:2:1: error: define-module: modules m1 use each other
e5.dyl:1:1: error: unknown module nosuch
e7.dyl:1:1: error: define-module: cannot export car, which is imported
e8b.dyl:5:5: error: define of f: the variable is imported from module a
\end{shell}

\noindent
A name that is used in a module but defined nowhere it can see draws a
warning, not an error, and an error only if the program actually reaches
it:

\begin{shell}
e1b.dyl:1:8: warning: g is not defined (module b)
$ ./e1
Dylan error: the variable g@b has no value
\end{shell}

\noindent
As the message shows, the compiler writes a variable as
\meta{name}\code{@}\meta{module} when the module matters. Here \code{g}
was defined in module \code{a} but not exported, so the \code{g} in module
\code{b} is a different, undefined variable.

\index{NAME-MODULE.c}
The generated C is organized by module. With \cmd{-k} (keep the C), a
program \code{app} with code in the modules \code{shapes} and \code{app}
leaves \file{app.c} (data, dispatch and \code{main}), \file{app.h},
\file{app-shapes.c} and \file{app-app.c}, one file of functions for each
module that has any, and \file{app-dylan.c}, the library's own functions.
Chapter~\ref{ch:dylanc} describes these files and the compiler's other
options.

\section{Modules in the interactive implementation}\label{sec:modules-lisp}

\index{interactive implementation!modules}
The Common Lisp implementation compiles one form at a time, so a module
must be declared before any code that uses it is compiled, and a module
must be declared after the modules it uses:

\begin{repl}
? (define-module m1 (use dylan) (use m2))

debugger invoked on a DYLAN::DYLAN-CONDITION ...
  Error: unknown module m2
\end{repl}

\noindent
Declaring the modules in order is enough. Redefining a module with
\code{define-module} replaces its clauses, and every module that uses it
is brought up to date at once, as \dylanbook{chapter 5} asks: variables
newly exported become visible immediately in the modules that import
them.

\subsection{At the read-eval-print loop}

\index{in-module@\code{dylan:in-module}}
The loop starts in \code{dylan-user}. Typing
\code{(dylan:in-module \meta{name})} at the \code{?} prompt switches the
module that later input is compiled in; from Lisp, the function
\code{(dylan:in-module "\meta{name}")}, given the name as a string, does
the same. The loop answers with
\code{;Module:} and the new module's name.

\begin{repl}
? (define-module counter (use dylan) (export next-count))
;Value: counter

? (dylan:in-module counter)
;Module: counter

? (define count 0)
;Value: count

? (define-method next-count () (set! count (+ count 1)) count)
;Value: next-count

? (dylan:in-module dylan-user)
;Module: dylan-user

? (define-module client (use dylan) (use counter))
;Value: client

? (dylan:in-module client)
;Module: client

? (next-count)
;Value: 1

? (next-count)
;Value: 2
\end{repl}

\noindent
Module \code{client} can call \code{next-count} but cannot see
\code{count}, which \code{counter} keeps to itself.
(Chapter~\ref{ch:lisp} describes the loop in full.)

\subsection{Files}

\code{dylan:compile-file} reads a file's header and compiles the file in
the module it names. Names are resolved to variables as the file is
compiled, so the compiled Lisp needs no module information to run, and a
compiled \code{define-module} form defines its module again when the
file is loaded. To
compile a program of several files, compile them in one Lisp session in
the order a reader would need them, the file of \code{define-module}
forms first, and then load the results in the same order:

\begin{shell}
$ sbcl --non-interactive --eval '(load "src/load-compiler.lisp")' \
    --eval '(dylan:compile-file "modules.dyl" "modules.lisp")' \
    --eval '(dylan:compile-file "shapes.dyl" "shapes.lisp")' \
    --eval '(dylan:compile-file "app.dyl" "app.lisp")'
$ sbcl --non-interactive --eval '(load "src/load-runtime.lisp")' \
    --eval '(load "modules.lisp")' --eval '(load "shapes.lisp")' \
    --eval '(load "app.lisp")'
\end{shell}

\noindent
Compiling \file{shapes.dyl} in a session that has not seen
\file{modules.dyl} fails with ``unknown module shapes''. When the Lisp
files are loaded, SBCL may print \code{WARNING: redefining ...} lines for
the small accessor functions that each file declares for the module
variables it shares with the others; they are harmless.

\section{Read-only variables}\label{sec:modules-readonly}

\index{read-only variable}
\dylanbook{Chapter 5} makes some module variables \term{read-only}: the
variable always holds the same object, though that object may itself be
changed. \code{define-class} and \code{define-generic-function} create
read-only variables, so a class or generic function cannot be replaced by
assignment. Variables made by \code{define} (or by \code{define-method}
alone) are ordinary and can be assigned with \code{set!}.

The two implementations check this at different times. The compiler to C
refuses the program:

\begin{shell}
e3.dyl:2:1: error: <p> holds a class or generic function and cannot be assigned
\end{shell}

\noindent
The interactive implementation signals an error when the assignment is
run:

\begin{repl}
? (define-class <p> (<object>))
;Value: <p>

? (set! <p> 3)

debugger invoked on a DYLAN::DYLAN-CONDITION ...
  Error: variable is read-only <p>
\end{repl}

\noindent
Methods can still be added to a read-only generic function: the variable
does not change, only the function in it.

\section{A complete multi-file program}\label{sec:modules-example}

This program has three files. The first declares the modules; it has no
header, so it is compiled in \code{dylan-user}.

\begin{dylan}
;;; modules.dyl
(define-module shapes
  (use dylan)
  (export <shape> <circle> <square> area describe-shape))

(define-module app
  (use dylan)
  (use shapes import: (<circle> <square> (describe-shape => show))))
\end{dylan}

\noindent
Module \code{shapes} exports five variables. Module \code{app} imports
only two classes and one function, and calls the function \code{show}.
It does not import \code{area} or \code{<shape>}.

\begin{dylan}
;;; shapes.dyl
module: shapes
author: A. Programmer

(define-class <shape> (<object>))
(define-class <circle> (<shape>)
  (radius init-keyword: radius:))
(define-class <square> (<shape>)
  (side init-keyword: side:))

(define-generic-function area (shape))
(define-method area ((c <circle>)) (* 3 (radius c) (radius c)))
(define-method area ((s <square>)) (* (side s) (side s)))

(define-method describe-shape ((s <shape>))
  (display "a shape of area ")
  (display (area s))
  (newline))
\end{dylan}

\noindent
The slot getters \code{radius} and \code{side} are owned by \code{shapes}
but not exported, so no other module can read the slots directly.

\begin{dylan}
;;; app.dyl
module: app

(show (make <circle> radius: 2))
(show (make <square> side: 3))
\end{dylan}

\noindent
With the compiler to C:

\begin{shell}
$ bin/dylanc -o app modules.dyl shapes.dyl app.dyl
$ ./app
a shape of area 12
a shape of area 9
\end{shell}

\noindent
With the interactive implementation, compile the three files and load
them as shown in Section~\ref{sec:modules-lisp}; the output is the same.

If \file{app.dyl} also called \code{(area (make <square> side: 1))}, the
compiler to C would warn that \code{area} is not defined in module
\code{app}, and the program would stop with an error on reaching the
call, because \code{app} did not import it.

A program meant to run in both implementations must keep to the names
both provide; Appendix~\ref{app:library} marks the ones that only one of
them has.
